\documentclass[11pt]{article}
\usepackage{mystyle}

\title{Rosen --- \S 8.1: Applications of Recurrence Relations\\ \large Selected Exercises}
\author{Alston}
\date{Semester 2, 2026}

\begin{document}
\maketitle

% ======================================================================
\begin{exercise}{11}
    \begin{enumerate}[label=(\alph*)]
        \item Find a recurrence relation for the number of ways to
        climb $n$ stairs if the person climbing the stairs can take
        one stair or two stairs at a time.
        \item What are the initial conditions?
        \item In how many ways can this person climb a flight of
        eight stairs?
    \end{enumerate}
\end{exercise}

% ======================================================================
\begin{exercise}{17}
    \textit{A string that contains only 0s, 1s, and 2s is called a
    ternary string.}
    \begin{enumerate}[label=(\alph*)]
        \item Find a recurrence relation for the number of ternary
        strings of length $n$ that do not contain consecutive symbols
        that are the same.
        \item What are the initial conditions?
        \item How many ternary strings of length six do not contain
        consecutive symbols that are the same?
    \end{enumerate}
\end{exercise}

% ======================================================================
\begin{exercise}{18}
    \begin{enumerate}[label=(\alph*)]
        \item Find a recurrence relation for the number of ternary
        strings of length $n$ that contain two consecutive symbols
        that are the same.
        \item What are the initial conditions?
        \item How many ternary strings of length six contain
        consecutive symbols that are the same?
    \end{enumerate}
\end{exercise}

% ======================================================================
\begin{exercise}{23}
    \begin{enumerate}[label=(\alph*)]
        \item Find the recurrence relation satisfied by $S_n$, where
        $S_n$ is the number of regions into which three-dimensional
        space is divided by $n$ planes if every three of the planes
        meet in one point, but no four of the planes go through the
        same point.
        \item Find $S_n$ using iteration.
    \end{enumerate}
\end{exercise}

% ======================================================================
\begin{exercise}{26}
    \begin{enumerate}[label=(\alph*)]
        \item Find a recurrence relation for the number of ways to
        completely cover a $2 \times n$ checkerboard with
        $1 \times 2$ dominoes.
        [\emph{Hint}: Consider separately the coverings where the
        position in the top right corner of the checkerboard is
        covered by a domino positioned horizontally and where it is
        covered by a domino positioned vertically.]
        \item What are the initial conditions for the recurrence
        relation in part (a)?
        \item How many ways are there to completely cover a
        $2 \times 17$ checkerboard with $1 \times 2$ dominoes?
    \end{enumerate}
\end{exercise}

% ======================================================================
\begin{exercise}{28}
    Show that the Fibonacci numbers satisfy the recurrence relation
    $f_n = 5f_{n-4} + 3f_{n-5}$ for $n = 5, 6, 7, \dots$, together with
    the initial conditions $f_0 = 0$, $f_1 = 1$, $f_2 = 1$, $f_3 = 2$,
    and $f_4 = 3$. Use this recurrence relation to show that $f_{5n}$
    is divisible by 5, for $n = 1, 2, 3, \dots$.
\end{exercise}

% ======================================================================
\begin{exercise}{29}
    Let $S(m,n)$ denote the number of onto functions from a set with
    $m$ elements to a set with $n$ elements. Show that $S(m,n)$
    satisfies the recurrence relation
    \[
    S(m,n) = n^m - \sum_{k=1}^{n-1} \binom{n}{k} S(m,k)
    \]
    whenever $m \geq n$ and $n > 1$, with the initial condition
    $S(m,1) = 1$.
\end{exercise}

% ======================================================================
\begin{exercise}{30}
    \textit{Recall (Example 5): $C_n$ denotes the number of ways to
    parenthesize a product of $n+1$ numbers so as to determine the
    order of multiplication; $C_0 = 1$ and
    $C_n = \sum_{k=0}^{n-1} C_k C_{n-1-k}$ for $n \geq 1$.}
    \begin{enumerate}[label=(\alph*)]
        \item Write out all the ways the product
        $x_0 \cdot x_1 \cdot x_2 \cdot x_3 \cdot x_4$ can be
        parenthesized to determine the order of multiplication.
        \item Use the recurrence relation developed in Example 5 to
        calculate $C_4$, the number of ways to parenthesize the
        product of five numbers so as to determine the order of
        multiplication. Verify that you listed the correct number of
        ways in part (a).
        \item Check your result in part (b) by finding $C_4$, using
        the closed formula for $C_n$ mentioned in the solution of
        Example 5.
    \end{enumerate}
\end{exercise}

% ======================================================================
\begin{exercise}{32}
    In the Tower of Hanoi puzzle, suppose our goal is to transfer all
    $n$ disks from peg 1 to peg 3, but we cannot move a disk directly
    between pegs 1 and 3. Each move of a disk must be a move involving
    peg 2. As usual, we cannot place a disk on top of a smaller disk.
    \begin{enumerate}[label=(\alph*)]
        \item Find a recurrence relation for the number of moves
        required to solve the puzzle for $n$ disks with this added
        restriction.
        \item Solve this recurrence relation to find a formula for
        the number of moves required to solve the puzzle for $n$
        disks.
        \item How many different arrangements are there of the $n$
        disks on three pegs so that no disk is on top of a smaller
        disk?
        \item Show that every allowable arrangement of the $n$ disks
        occurs in the solution of this variation of the puzzle.
    \end{enumerate}
\end{exercise}

% ======================================================================
\begin{exercise}{35}
    \textit{Recall the Josephus problem: $n$ people, numbered 1 to
    $n$, stand around a circle. In each stage, every second person
    still left alive is eliminated until only one survives. Denote
    the number of the survivor by $J(n)$.}

    Show that $J(n)$ satisfies the recurrence relation
    $J(2n) = 2J(n) - 1$ and $J(2n+1) = 2J(n) + 1$, for $n \geq 1$, and
    $J(1) = 1$.
\end{exercise}

% ======================================================================
\begin{exercise}{41}
    \textit{Recall the Reve's puzzle: the variation of the Tower of
    Hanoi puzzle with four pegs and $n$ disks. The Frame--Stewart
    algorithm moves the disks from peg 1 to peg 4 by choosing an
    integer $k$ with $1 \leq k \leq n$, recursively moving the $n-k$
    smallest disks to peg 2 (using all four pegs), then moving the
    $k$ largest disks to peg 4 (using the three-peg algorithm, without
    the peg holding the $n-k$ smallest disks), and finally recursively
    moving the $n-k$ smallest disks to peg 4.}

    Show that if $R(n)$ is the number of moves used by the
    Frame--Stewart algorithm to solve the Reve's puzzle with $n$
    disks, where $k$ is chosen to be the smallest integer with
    $n \leq k(k+1)/2$, then $R(n)$ satisfies the recurrence relation
    $R(n) = 2R(n-k) + 2k - 1$, with $R(0) = 0$ and $R(1) = 1$.
\end{exercise}

% ======================================================================
\begin{exercise}{45}
    Show that $R(n)$ (the number of moves used by the Frame--Stewart
    algorithm, from Exercise 41) is $O(\sqrt{n}\, 2^{\sqrt{n}})$.
\end{exercise}

% ======================================================================
\begin{exercise}{52}
    \textit{Recall: the backward differences of a sequence $\{a_n\}$
    are defined recursively by $\nabla a_n = a_n - a_{n-1}$ and
    $\nabla^{k+1} a_n = \nabla^k a_n - \nabla^k a_{n-1}$.}

    Show that any recurrence relation for the sequence $\{a_n\}$ can
    be written in terms of $a_n, \nabla a_n, \nabla^2 a_n, \dots$. The
    resulting equation involving the sequence and its differences is
    called a \emph{difference equation}.
\end{exercise}

% ======================================================================
\begin{exercise}{56}
    In this exercise we will develop a dynamic programming algorithm
    for finding the maximum sum of consecutive terms of a sequence of
    real numbers. That is, given a sequence of real numbers
    $a_1, a_2, \dots, a_n$, the algorithm computes the maximum sum
    $\sum_{i=j}^{k} a_i$ where $1 \leq j \leq k \leq n$.
    \begin{enumerate}[label=(\alph*)]
        \item Show that if all terms of the sequence are nonnegative,
        this problem is solved by taking the sum of all terms. Then,
        give an example where the maximum sum of consecutive terms is
        not the sum of all terms.
        \item Let $M(k)$ be the maximum of the sums of consecutive
        terms of the sequence ending at $a_k$. That is,
        $M(k) = \max_{1 \leq j \leq k} \sum_{i=j}^{k} a_i$. Explain
        why the recurrence relation $M(k) = \max(M(k-1)+a_k, a_k)$
        holds for $k = 2, \dots, n$.
        \item Use part (b) to develop a dynamic programming algorithm
        for solving this problem.
        \item Show each step your algorithm from part (c) uses to
        find the maximum sum of consecutive terms of the sequence
        $2, -3, 4, 1, -2, 3$.
        \item Show that the worst-case complexity in terms of the
        number of additions and comparisons of your algorithm from
        part (c) is linear.
    \end{enumerate}
\end{exercise}

% ======================================================================
\begin{exercise}{57}
    Dynamic programming can be used to develop an algorithm for
    solving the matrix-chain multiplication problem. This is the
    problem of determining how the product $A_1 A_2 \cdots A_n$ can be
    computed using the fewest integer multiplications, where
    $A_1, A_2, \dots, A_n$ are $m_1 \times m_2$, $m_2 \times m_3$,
    $\dots$, $m_n \times m_{n+1}$ matrices, respectively, and each
    matrix has integer entries. Recall that by the associative law,
    the product does not depend on the order in which the matrices
    are multiplied.
    \begin{enumerate}[label=(\alph*)]
        \item Show that the brute-force method of determining the
        minimum number of integer multiplications needed to solve a
        matrix-chain multiplication problem has exponential
        worst-case complexity.
        [\emph{Hint}: Do this by first showing that the order of
        multiplication of matrices is specified by parenthesizing the
        product, then use the Catalan-number count of parenthesizations
        (Exercise 30) together with its exponential lower bound.]
        \item Denote by $A_{ij}$ the product $A_i A_{i+1} \cdots A_j$,
        and $M(i,j)$ the minimum number of integer multiplications
        required to find $A_{ij}$. Show that if the least number of
        integer multiplications are used to compute $A_{ij}$, where
        $i < j$, by splitting the product into the product of $A_i$
        through $A_k$ and the product of $A_{k+1}$ through $A_j$,
        then the first $k$ terms must be parenthesized so that
        $A_{i,k}$ is computed in the optimal way using $M(i,k)$
        integer multiplications and $A_{k+1,j}$ must be parenthesized
        so that $A_{k+1,j}$ is computed in the optimal way using
        $M(k+1,j)$ integer multiplications.
        \item Explain why part (b) leads to the recurrence relation
        \[
        M(i,j) = \min_{i \leq k < j} \big(M(i,k) + M(k+1,j) + m_i
        m_{k+1} m_{j+1}\big) \quad \text{if } 1 \leq i < j \leq n.
        \]
        \item Use the recurrence relation in part (c) to construct an
        efficient algorithm for determining the order the $n$
        matrices should be multiplied to use the minimum number of
        integer multiplications. Store the partial results $M(i,j)$
        as you find them so that your algorithm will not have
        exponential complexity.
        \item Show that your algorithm from part (d) has $O(n^3)$
        worst-case complexity in terms of multiplications of
        integers.
    \end{enumerate}
\end{exercise}

\end{document}